\subsection{Image of a measure}

DEFINITION 4. --- Let \(\pi\) be a mapping of T into a topological space X. One says that \(\pi\) is \(\mu\) -proper if \(\pi\) is \(\mu\) -measurable, and if every point \(x\) of X admits a neighborhood V such that \(\mu^{\bullet}\left(\overline{\pi}^{-1}(\mathrm{V})\right) < +\infty\) .

Remarks. -- 1) When T and X are locally compact, this definition is equivalent to that of Ch. V, §6, No.~1.

\begin{enumerate}
\def\labelenumi{\arabic{enumi})}
\setcounter{enumi}{1}
\item
  A proper continuous mapping (GT, I, §10, No.~1, Def. 1) of T into X is \(\mu\) -proper for every measure \(\mu\) . For, let \(x \in X\) ; since \(\overline{\pi}^{-1}(x)\) is compact (loc. cit., No.~2, Th. 1), the set \(\overline{\pi}^{-1}(x)\) has an open neighborhood H such that \(\mu^{\bullet}(\mathrm{H}) < +\infty\) . Set \(\mathrm{V} = \mathrm{X} - \pi(\mathrm{T} - \mathrm{H})\) ; since \(\pi\) is closed, V is open in X, contains \(x\) , and satisfies \(\overline{\pi}^{-1}(\mathrm{V}) \subset \mathrm{H}\) , whence \(\mu^{\bullet}(\overline{\pi}^{-1}(\mathrm{V})) \leqslant \mu^{\bullet}(\mathrm{H}) < +\infty\) .
\item
  If \(\mu\) is bounded, every \(\mu\) -measurable mapping of T into X is \(\mu\) -proper.
\item
  If \(\theta\) is a complex measure on \(\mathbf{T}\) , \(\pi\) is said to be \(\theta\) -proper if \(\pi\) is proper for the positive measure \(|\theta|\) .
\end{enumerate}

PROPOSITION 4. --- Let \(\pi\) be a \(\mu\) -proper mapping of T into a topological space X. There exists one and only one measure \(\nu\) on X such that \(\nu^{\bullet}\) is equal to the image encumbrance \(\pi(\mu^{\bullet})\) (§1, No.~1), in other words, such that \(\nu^{\bullet}(g) = \mu^{\bullet}(g \circ \pi)\) for all \(g \in \mathcal{F}_{+}(\mathbf{X})\) .

Uniqueness is obvious (§1, No.~2, Cor. of Prop. 2). To establish existence, we shall first treat the case that \(\mu\) is carried by a compact set K, such that the restriction of \(\pi\) to K is continuous. Then \(L = \pi(K)\) is compact; let \(\pi'\) be the continuous mapping of K into L induced by \(\pi\) , and let \(\nu'\) be the image measure \(\pi'(\mu_K)\) on L, \(\nu\) the measure on X defined by \(\nu'\) (§1, No.~3, Example 2). For every \(g \in \mathcal{F}_+(X)\) ,

\[
\nu^ {\bullet} (g) = \nu^ {\prime \bullet} (g _ {\mathrm{L}}) = \mu_ {\mathrm{K}} ^ {\bullet} (g _ {\mathrm{L}} \circ \pi^ {\prime}) = \mu_ {\mathrm{K}} ^ {\bullet} ((g \circ \pi) _ {\mathrm{K}}) = \mu^ {\bullet} ((g \circ \pi) _ {\mathrm{K}} ^ {0}) = \mu^ {\bullet} (g \circ \pi)
\]

(we have used successively formula (3) of §1, No.~3; Prop. 2 of Ch. V, §6, No.~2; the definition of \(\mu_{\mathrm{K}}^{\bullet}\) ; and the fact that \(\mu\) is carried by K). In other words, \(\nu^{\bullet} = \pi(\mu^{\bullet})\) .

Let us now pass to the general case; by Props. 10 and 9 of §1, No.~8, \(\mu\) is the sum of a summable family \((\mu_{\alpha})_{\alpha \in \mathbf{A}}\) of measures with compact support, such that the restriction of \(\pi\) to the support \(\mathbf{K}_{\alpha}\) of \(\mu_{\alpha}\) is continuous for every \(\alpha \in \mathbf{A}\) . The special case treated above permits associating to each measure \(\mu_{\alpha}\) on \(\mathbf{T}\) a measure \(\nu_{\alpha}\) on \(\mathbf{X}\) such that \(\nu_{\alpha}^{\bullet} = \pi (\mu_{\alpha}^{\bullet})\) . Then, for \(g \in \mathcal{F}_{+}(\mathbf{X})\) ,

\[
\sum_ {\alpha \in A} \nu_ {\alpha} ^ {\bullet} (g) = \sum_ {\alpha \in A} \mu_ {\alpha} ^ {\bullet} (g \circ \pi) = \mu^ {\bullet} (g \circ \pi).
\]

The encumbrance \(\pi (\mu^{\bullet})\) is locally bounded, since \(\pi\) is \(\mu\) -proper; the family \((\nu_{\alpha})_{\alpha \in \mathbf{A}}\) is therefore summable (§1, No.~7, Prop. 7), and its sum \(\nu\) satisfies the statement.

DEFINITION 5. --- If \(\pi\) is a \(\mu\) -proper mapping of T into a topological space X, the unique measure \(\nu\) on X such that \(\nu^{\bullet}(g) = \mu^{\bullet}(g \circ \pi)\) for all \(g \in \mathcal{F}_{+}(\mathrm{X})\) is called the image measure of \(\mu\) under \(\pi\) , and is denoted \(\pi(\mu)\) .

Example. --- Let K be a compact subspace of T, i the canonical injection of K into T, and \(\lambda\) a measure on K; since i is continuous and \(\lambda\) is bounded, i is \(\lambda\) -proper. Formula (3) of §1, No.~3 shows that the ``measure on T defined by \(\lambda\) '' is the image measure \(i(\lambda)\) .

Remark 5). --- If \(\theta\) is a real measure and \(\pi\) is \(\theta\) -proper, then \(\pi\) is proper for the measures \(\theta^{+}\) and \(\theta^{-}\) ; one then sets \(\pi(\theta) = \pi(\theta^{+}) - \pi(\theta^{-})\) . If \(\theta\) is a complex measure and \(\pi\) is \(\theta\) -proper, then \(\pi\) is proper for the real measures \(\mathcal{R}(\theta)\) and \(\mathcal{I}(\theta)\) ; one then sets

\[
\pi (\theta) = \pi (\mathcal {R} (\theta)) + i \pi (\mathcal {I} (\theta)).
\]

PROPOSITION 5. --- Let \(\pi\) be a \(\mu\) -proper mapping of T into a topological space X, and let \(f\) be a mapping of X into a topological space F (Hausdorff or not). For \(f\) to be \(\pi(\mu)\) -measurable, it is necessary and sufficient that \(f \circ \pi\) be \(\mu\) -measurable.

Let us take up again the proof of Prop. 4, and commence with the special case treated at the beginning, with the same notations; \(g\) is measurable for the measure \(\pi(\mu) = \nu\) if and only if \(g_{\mathrm{L}}\) is \(\nu'\) -measurable (§1, No.~5, Example); now, this is equivalent to saying that \(g_{\mathrm{L}} \circ \pi' = (g \circ \pi)_{\mathrm{K}}\) is \(\mu_{\mathrm{K}}\) -measurable (Ch. V, §6, No.~2, Prop. 3), and finally that \(g \circ \pi\) is \(\mu\) -measurable (§1, No.~5, Example). Let us now pass to the general case, with the same notations as in the proof of Prop. 4; \(f\) is \(\nu\) -measurable if and only if \(f\) is \(\nu_{\alpha}\) -measurable for every \(\alpha \in \mathrm{A}\) (§1, No.~7, Prop. 8), hence if and only if \(f \circ \pi\) is \(\mu_{\alpha}\) -measurable for every \(\alpha \in \mathbf{A}\) (the preceding special case) and finally if and only if \(f \circ \pi\) is \(\mu\) -measurable (§1, No.~7, Prop. 8).

COROLLARY. --- Let X and Y be two topological spaces, \(\pi\) a \(\mu\) -proper mapping of T into X, and \(\pi'\) a \(\pi(\mu)\) -proper mapping of X into Y. The mapping \(\pi'' = \pi' \circ \pi\) is then \(\mu\) -proper, and \(\pi''(\mu) = \pi'(\pi(\mu))\) (`transitivity of images of measures').

For, \(\pi''\) is \(\mu\) -measurable (Prop. 5). Set \(\mu' = \pi(\mu)\) ; the image encumbrance \(\pi'(\mu'^{\bullet}) = \pi'(\pi(\mu^{\bullet}))\) is obviously equal to \(\pi''(\mu^{\bullet})\) . Since it is locally bounded, \(\pi''\) is \(\mu\) -proper. The measures \(\pi''(\mu)\) and \(\pi'(\mu')\) then have the same essential upper integral, hence are equal.

PROPOSITION 6. --- Let \(\pi\) be a \(\mu\) -proper mapping of T into a topological space X, and let B be a \(\pi(\mu)\) -measurable subset of X. Set A = \(\overline{\pi}^{-1}(B)\) , and denote by \(\pi'\) the mapping of A into B that coincides with \(\pi\) on A. The set A is then \(\mu\) -measurable, \(\pi_A\) and \(\pi'\) are \(\mu_A\) -proper, and

\[
\left(\pi (\mu)\right) _ {\mathrm{B}} = \left(\pi_ {\mathrm{A}} (\mu_ {\mathrm{A}})\right) _ {\mathrm{B}} = \pi^ {\prime} (\mu_ {\mathrm{A}}).\tag{2}
\]

The set A is \(\mu\) -measurable by Prop. 5 applied to \(\varphi_{B}\) ; the mapping \(\pi_{A}\) is clearly \(\mu_{A}\) -measurable by the definition of induced measures (No.~1), and it follows that \(\pi'\) is measurable. Let f be an element of \(\mathcal{F}_{+}(B)\) ; denoting by zero exponents the extensions by 0 in X and in T, we have

\[
\left(\pi (\mu) _ {\mathrm{B}}\right) ^ {\bullet} (f) = \pi (\mu) ^ {\bullet} \left(f ^ {0}\right) = \mu^ {\bullet} \left(f ^ {0} \circ \pi\right) = \mu^ {\bullet} \left(\left(f \circ \pi^ {\prime}\right) ^ {0}\right) = \mu_ {\mathrm{A}} ^ {\bullet} \left(f \circ \pi^ {\prime}\right),
\]

whence \(\left(\pi(\mu)_{\mathrm{B}}\right)^{\bullet} = \pi'(\mu_{\mathrm{A}}^{\bullet})\) . Since the encumbrance \(\left(\pi(\mu)_{\mathrm{B}}\right)^{\bullet}\) is locally bounded, the same is true of \(\pi'(\mu_{\mathrm{A}}^{\bullet})\) and therefore \(\pi'\) is \(\mu_{A}\) -proper. The measures \(\pi'(\mu_{\mathrm{A}})\) and \(\left(\pi(\mu)\right)_{\mathrm{B}}\) have the same essential upper integral, hence are equal. The other relation may be established in an analogous manner.

PROPOSITION 7. --- Let T be a subspace of a topological space X, and let i be the injection of T into X.

\begin{enumerate}
\def\labelenumi{\alph{enumi})}
\item
  If \(\mu\) is a measure on T, and if \(i\) is \(\mu\) -proper, then the measure \(i(\mu)\) is concentrated on T, and one has \((i(\mu))_{T} = \mu\) .
\item
  If \(\lambda\) is a measure on \(X\) such that \(T\) is \(\lambda\) -measurable, then \(i\) is \(\lambda_{T}\) -proper, and \(i(\lambda_T) = \varphi_T \cdot \lambda\) .
\item
  Set \(\nu = i(\mu)\) ; the relation \(\nu^{\bullet}(\mathbf{A}) = \mu^{\bullet}(\mathbf{A} \cap \mathbf{T})\) , applied to \(\mathbf{A} = \mathbf{X} - \mathbf{T}\) , shows that \(\nu\) is concentrated on \(\mathbf{T}\) . The relation \(\nu_{\mathrm{T}} = \mu\) is a special case of the relation (2), on taking \(\mathbf{B} = \mathbf{T} = \mathbf{A}\) .
\item
  Let \(f\) be a positive function defined on \(\mathbf{X}\) ; setting \(\mu = \lambda_{\mathrm{T}}\) , one has \(\mu^{\bullet}(f \circ i) = \lambda_{\mathrm{T}}^{\bullet}(f_{\mathrm{T}}) = \lambda^{\bullet}(f\varphi_{\mathrm{T}}) \leqslant \lambda^{\bullet}(f)\) (Prop. 1); it follows that \(i\) is \(\mu\) -proper. On the other hand, \(\mu^{\bullet}(f \circ i)\) (resp. \(\lambda^{\bullet}(f\varphi_{\mathrm{T}})\) ) is the essential upper integral of f with respect to \(i(\mu)\) (resp. \(\varphi_{T} \cdot \lambda\) ). These two measures are therefore equal.
\end{enumerate}

Remark 6). --- Let \(\pi\) be a \(\mu\) -proper mapping of T into a topological space X. One reduces the theory of integration with respect to the image measure \(\nu = \pi(\mu)\) to the theory treated in Ch. V, §6, in the following way. Let \((\mathrm{K}_{\alpha})_{\alpha \in \mathrm{A}}\) (resp. \((\mathrm{L}_{\beta})_{\beta \in \mathrm{B}}\) ) be a crushing of T (resp. of X) for \(\mu\) (resp. for \(\nu\) ), and set \(N = T - \bigcup_{\alpha \in A} K_{\alpha}\) , \(P = X - \bigcup_{\beta \in B} L_{\beta}\) . We can suppose that the restriction of \(\pi\) to each \(K_{\alpha}\) is continuous (§1, No.~8, Prop. 10). Let \(T', X'\) be the locally compact spaces constructed as in the Scholium of §1, No.~8 and let \(\mu'\) and \(\nu'\) be the measures on these spaces associated with \(\mu\) and \(\nu\) . The topology of \(T'\) being the sum of the topologies of the subspaces \(K_{\alpha}\) and the discrete topology on N, \(\pi\) is a continuous mapping of \(T'\) into X and the relation \(\mu'^{\bullet}(g \circ \pi) = \mu^{\bullet}(g \circ \pi) = \nu^{\bullet}(g)\) (for \(g \in \mathcal{F}_{+}(X)\) ) shows that \(\pi\) is \(\mu'\) -proper and that \(\pi(\mu') = \nu\) . On the other hand, the identity mapping \(i\) of X onto \(X'\) is \(\nu\) -proper, and \(i(\nu) = \nu'\) . It follows that \(\pi\) is a \(\mu'\) -proper mapping of \(T'\) into \(X'\) , and that the image of \(\mu'\) under \(\pi\) is \(\nu'\) (Cor. of Prop. 5). We leave to the reader the task of transcribing the results of Ch. V, §6.

